\documentclass{article}
\usepackage{amsmath, amssymb, amsthm}
\usepackage{hyperref}
\usepackage{geometry}
\geometry{margin=1in}

\title{Erd\H{o}s Problem \#213}
\date{Last updated: 2025-08-31}
\author{Source: erdosproblems.com}

\begin{document}
\maketitle

\noindent\textbf{Status:} OPEN \quad \textbf{No prize}

\noindent\textbf{Tags:} geometry, distances

\medskip
\noindent\textbf{Problem Statement:}

Let $n\geq 4$. Are there $n$ points in $\mathbb{R}^2$, no three on a line and no four on a circle, such that all pairwise distances are integers?

\bigskip
\noindent\textbf{Remarks:}

Anning and Erd\H{o}s \cite{AnEr45} proved there cannot exist an infinite such set. Harborth constructed such a set when $n=5$. The best construction to date, due to Kreisel and Kurz \cite{KK08}, has $n=7$. 

Ascher, Braune, and Turchet \cite{ABT20} have shown that there is a uniform upper bound on the size of such a set, conditional on the Bombieri-Lang conjecture. Greenfeld, Iliopoulou, and Peluse \cite{GIP24} have shown (unconditionally) that any such set must be very sparse, in that if $S\subseteq [-N,N]^2$ has no three on a line and no four on a circle, and all pairwise distances integers, then\[\lvert S\rvert \ll (\log N)^{O(1)}.\]See also [130].

\bigskip
\noindent\textbf{References:}

[ABT20] Ascher, K. and Braune, L. and Turchet, A., The Erd\H{o}s-Ulam problem, Lang's conjecture, and uniformity. arXiv:1901.02616 (2020).
 
 [AnEr45] Anning, Norman H. and Erd\H{o}s, Paul, Integral distances. Bull. Amer. Math. Soc. (1945), 598-600.
 
 [GIP24] Greenfeld, R. and Iliopoulou, M. and Peluse, S., On integer distance sets. arXiv:2401.10821 (2024).
 
 [KK08] Kreisel, Tobias and Kurz, Sascha, There are integral heptagons, no three points on a line, on four on a circle. Discrete Comput. Geom. (2008), 786-790.

\bigskip
\noindent\small{Source: \url{https://www.erdosproblems.com/213}}
\end{document}
